\hypertarget{python-1.x-the-pyobject-model-reference-counting-core}{%
\chapter{Python 1.x: The PyObject Model \& Reference Counting
Core}\label{python-1.x-the-pyobject-model-reference-counting-core}}

\hypertarget{the-unified-pyobject-pyvarobject-structs}{%
\subsection{\texorpdfstring{2.1 The Unified \texttt{PyObject} \&
\texttt{PyVarObject}
Structs}{2.1 The Unified PyObject \& PyVarObject Structs}}\label{the-unified-pyobject-pyvarobject-structs}}

In CPython, \textbf{every python value is a heap-allocated C struct}.
There are no primitive variables on the stack. The fundamental base
structure is defined in \passthrough{\lstinline!Include/object.h!}:

\begin{lstlisting}[language=C]
/* CPython source definition from Include/object.h */
#ifdef Py_TRACE_REFS
/* Doubly-linked list pointers for active heap tracing in debug builds */
#define _PyObject_HEAD_EXTRA            \
    struct _object *_ob_next;           \
    struct _object *_ob_prev;
#else
#define _PyObject_HEAD_EXTRA
#endif

typedef struct _object {
    _PyObject_HEAD_EXTRA                /* Doubly-linked list pointers */
    Py_ssize_t ob_refcnt;               /* Reference count counter */
    struct _typeobject *ob_type;        /* Pointer to the type definition object */
} PyObject;

typedef struct {
    PyObject ob_base;                   /* Base type fields */
    Py_ssize_t ob_size;                 /* Size of the variable portion (e.g. sequence length) */
} PyVarObject;
\end{lstlisting}

\hypertarget{comparison-to-c-memory-layout}{%
\subsubsection{1. Comparison to C++ Memory
Layout}\label{comparison-to-c-memory-layout}}

Unlike C++, where objects are laid out contiguously on the stack or heap
based on their static type declarations and resolved via virtual
function tables (\passthrough{\lstinline!vtable!}), CPython objects are
entirely dynamic: 1. \textbf{Uniform Pointer Width}: Every variable in
Python is a pointer (\passthrough{\lstinline!PyObject*!}), consuming
exactly 8 bytes (on 64-bit platforms). 2. \textbf{No VTable
Indirection}: Polymorphism is resolved dynamically by dereferencing the
type pointer \passthrough{\lstinline!ob\_type!} at runtime to lookup
operations.

\begin{lstlisting}
Python Reference Pointer:
[Name: x] ---> [ Heap Allocation: PyObject ]
               +--------------------------------------+
               | _PyObject_HEAD_EXTRA (Tracking)      |
               | ob_refcnt: 1                         |
               | ob_type: ----> [ PyTypeObject (int)] |
               | Integer Data Value Payload           |
               +--------------------------------------+
\end{lstlisting}

\begin{center}\rule{0.5\linewidth}{0.5pt}\end{center}

\hypertarget{the-type-object-slot-system}{%
\subsection{2.2 The Type-Object Slot
System}\label{the-type-object-slot-system}}

The type of an object is defined by the
\passthrough{\lstinline!PyTypeObject!} struct (which is itself a
subclass of \passthrough{\lstinline!PyObject!}).

\hypertarget{type-struct-slots-layout}{%
\subsubsection{1. Type Struct Slots
layout}\label{type-struct-slots-layout}}

The type struct defines a series of ``slots'' (pointers to C functions)
that determine how the object behaves when subjected to operators.

\begin{lstlisting}[language=C]
/* Conceptual CPython definition from Include/cpython/object.h */
struct _typeobject {
    PyObject_VAR_HEAD
    const char *tp_name;                 /* For printing, in format <module>.<name> */
    Py_ssize_t tp_basicsize, tp_itemsize; /* For allocation */
    
    /* Type Slots (Static Function Pointers) */
    destructor tp_dealloc;               /* Deallocation handler */
    reprfunc tp_repr;                    /* Representation builder */
    getattrfunc tp_getattr;              /* Attribute lookup hook */
    setattrfunc tp_setattr;              /* Attribute assignment hook */
    hashfunc tp_hash;                    /* Hashing calculation function pointer */
    
    /* Protocol Table Slots */
    PyNumberMethods *tp_as_number;       /* Math function pointers (e.g. nb_add) */
    PySequenceMethods *tp_as_sequence;   /* Indexing function pointers (e.g. sq_item) */
    PyMappingMethods *tp_as_mapping;     /* Map lookup function pointers (e.g. mp_subscript) */
};
typedef struct _typeobject PyTypeObject;
\end{lstlisting}

\hypertarget{dunder-methods-to-slot-mapping}{%
\subsubsection{2. Dunder Methods to Slot
Mapping}\label{dunder-methods-to-slot-mapping}}

When defining a class in Python, the interpreter automatically maps
custom dunder methods to these internal C slots: *
\passthrough{\lstinline!\_\_init\_\_!} maps to
\passthrough{\lstinline!tp\_init!}. *
\passthrough{\lstinline!\_\_new\_\_!} maps to
\passthrough{\lstinline!tp\_new!}. *
\passthrough{\lstinline!\_\_hash\_\_!} maps to
\passthrough{\lstinline!tp\_hash!}. *
\passthrough{\lstinline!\_\_add\_\_!} maps to
\passthrough{\lstinline!tp\_as\_number->nb\_add!}. *
\passthrough{\lstinline!\_\_getitem\_\_!} maps to
\passthrough{\lstinline!tp\_as\_mapping->mp\_subscript!} or
\passthrough{\lstinline!tp\_as\_sequence->sq\_item!}.

\begin{lstlisting}[language=Python]
# Customizing slots via class definitions
class CustomObject:
    def __init__(self, val):
        self.val = val

    def __hash__(self):
        # Maps to the tp_hash C slot
        return hash(self.val)

obj = CustomObject("test")
print("Hash via type slot:", hash(obj))
\end{lstlisting}

\begin{center}\rule{0.5\linewidth}{0.5pt}\end{center}

\hypertarget{reference-counting-lifecycle}{%
\subsection{2.3 Reference Counting
Lifecycle}\label{reference-counting-lifecycle}}

CPython manages memory directly via reference counting. Reference counts
are modified using C macros.

\hypertarget{c-level-reference-modification-macros}{%
\subsubsection{1. C-Level Reference Modification
Macros}\label{c-level-reference-modification-macros}}

Defined in \passthrough{\lstinline!object.h!}, the macros expand as:

\begin{lstlisting}[language=C]
#define Py_INCREF(op) (                         \
    _Py_INC_REFTOTAL  _Py_REF_DEBUG_EXTRA      \
    (op)->ob_refcnt++)

#define Py_DECREF(op) do {                      \
    if (_Py_DEC_REFTOTAL  _Py_REF_DEBUG_EXTRA   \
        --(op)->ob_refcnt == 0)                 \
        _Py_Dealloc((PyObject *)(op));          \
} while (0)
\end{lstlisting}

To prevent segmentation faults when referencing a null pointer, CPython
provides the NULL-safe macros \passthrough{\lstinline!Py\_XINCREF(op)!}
and \passthrough{\lstinline!Py\_XDECREF(op)!}, which check
\passthrough{\lstinline"if (op != NULL)"} before modifying the counts.

\hypertarget{the-deallocation-pipeline}{%
\subsubsection{2. The Deallocation
Pipeline}\label{the-deallocation-pipeline}}

When \passthrough{\lstinline!Py\_DECREF(obj)!} drops the reference count
to zero: 1. \textbf{Invoke Deallocator}: The VM invokes
\passthrough{\lstinline!obj->ob\_type->tp\_dealloc(obj)!}. 2.
\textbf{Recursive Decrements}: The deallocator function decrements
references of any nested objects inside this object. For example, if a
list is deallocated, it decrements the reference counts of all items it
contains. 3. \textbf{Release Memory}: The deallocator releases the raw
memory back to PyMalloc or the system allocator.

\hypertarget{leak-tracking-in-c-extensions}{%
\subsubsection{3. Leak Tracking in C
Extensions}\label{leak-tracking-in-c-extensions}}

In custom C extensions, failing to decrement references leads to memory
leaks. In debug builds, compiling CPython with
\passthrough{\lstinline!Py\_TRACE\_REFS!} tracks all active references
in a global doubly-linked list. We can trace reference counts using the
\passthrough{\lstinline!sys!} module:

\begin{lstlisting}[language=Python]
import sys

# Reference counts tracking
my_list = [100, 200]
print("Starting references count:", sys.getrefcount(my_list) - 1)

ref_a = my_list
ref_b = my_list
print("References after bindings:", sys.getrefcount(my_list) - 1)

del ref_a
print("References after deleting one binding:", sys.getrefcount(my_list) - 1)
\end{lstlisting}

\begin{center}\rule{0.5\linewidth}{0.5pt}\end{center}

\hypertarget{built-in-caching-and-interning-optimizations}{%
\subsection{2.4 Built-in Caching and Interning
Optimizations}\label{built-in-caching-and-interning-optimizations}}

To avoid allocating memory for common objects, CPython uses global
caches:

\hypertarget{small-integer-cache}{%
\subsubsection{1. Small Integer Cache}\label{small-integer-cache}}

CPython pre-allocates an array of integer objects from
\passthrough{\lstinline!-5!} to \passthrough{\lstinline!256!}
(\passthrough{\lstinline!NSMALLNEGINTS!} and
\passthrough{\lstinline!NSMALLPOSINTS!}). * \textbf{Mechanism}: Any
assignment in this range returns a pointer to the pre-existing static
object, bypassing heap allocations. * \textbf{C-level array}: Inside
\passthrough{\lstinline!objects/longobject.c!}, this is defined as:

\begin{lstlisting}[language=C]
static PyLongObject small_ints[NSMALLNEGINTS + NSMALLPOSINTS];
\end{lstlisting}

\hypertarget{string-interning}{%
\subsubsection{2. String Interning}\label{string-interning}}

String literals resembling identifiers are automatically
\textbf{interned}. * \textbf{Mechanism}: Interned strings are stored in
a global dictionary inside \passthrough{\lstinline!unicodeobject.c!}. If
a string is already interned, any new instantiation returns the existing
string pointer. * \textbf{Performance}: This allows matching strings via
fast pointer comparison (\(O(1)\) address checks) instead of
byte-by-byte comparison (\(O(N)\) string scans).

\begin{lstlisting}[language=C]
/* Under the hood string comparison */
if (str1 == str2) {
    return 1; /* Match! (O(1) address check) */
}
\end{lstlisting}

\hypertarget{empty-collections-and-singletons}{%
\subsubsection{3. Empty Collections and
Singletons}\label{empty-collections-and-singletons}}

To save memory: * \textbf{Empty Tuple Cache}: CPython allocates exactly
one global empty tuple object
(\passthrough{\lstinline!\&\_Py\_EmptyTuple!}). Every call to
\passthrough{\lstinline!tuple()!} returns a pointer to this singleton. *
\textbf{Built-in Singletons}: \passthrough{\lstinline!None!},
\passthrough{\lstinline!True!}, \passthrough{\lstinline!False!},
\passthrough{\lstinline!Ellipsis!}, and
\passthrough{\lstinline!NotImplemented!} are statically allocated global
structures.

\begin{lstlisting}[language=Python]
# Small Integer Caching Check
int_a = 256
int_b = 256
print("Are 256 pointers identical?", int_a is int_b) # True (cached)

int_c = 257
int_d = 257
print("Are 257 pointers identical?", int_c is int_d) # False (dynamically allocated)

# String Interning Check
str_a = "godhood_string"
str_b = "godhood_string"
print("Are literal strings interned?", str_a is str_b) # True

# Empty tuple check
tup_a = ()
tup_b = tuple()
print("Are empty tuples cached singletons?", tup_a is tup_b) # True
\end{lstlisting}

\begin{center}\rule{0.5\linewidth}{0.5pt}\end{center}
